\begin{figure}[ht]
\centering
\newcommand{\pipelineheading}{\fontsize{8.85}{10.0}\selectfont}
\newcommand{\pipelinesupport}{\fontsize{6.46}{7.5}\selectfont}
\begin{tikzpicture}[
  font=\pipelineheading,
  >=Latex,
  box/.style={draw=black!55, rounded corners=2pt, align=center,
              inner sep=5pt, minimum height=9mm, fill=black!3},
  wide/.style={box, text width=74mm},
  stage/.style={box, fill=blue!8, draw=blue!45, text width=30mm},
  output/.style={box, fill=orange!10, draw=orange!55, text width=30mm},
  fixed/.style={box, fill=green!8, draw=green!45, text width=24mm}
]
\node[wide] (record) at (1.3,3.3) {\textbf{Shared record}\\
  {\pipelinesupport question, four answer choices, and the fixed description in the description arms}};
\node[wide] (constructor) at (1.3,1.7) {\textbf{Image condition}\\
  {\pipelinesupport source, grey, mismatch, mask, noise, mirror}};

\node[stage] (stage1) at (-2.9,0.0) {\textbf{Stage~1}\\
  {\pipelinesupport own adaptation}};
\node[stage] (stage2) at (3.2,0.0) {\textbf{Stage~2}\\
  {\pipelinesupport own adaptation}};
\node[fixed] (gold) at (7.0,0.0) {gold answer\\
  {\pipelinesupport supplied, unchanged}};

\node[output] (answer) at (-2.9,-1.8) {predicted answer};
\node[output] (rationale) at (3.2,-1.8) {generated rationale};

\draw[->,thick] (record) -- (constructor);
\draw[->,thick] (constructor.south) -- ++(0,-0.35) -| (stage1.north);
\draw[->,thick] (constructor.south) -- ++(0,-0.35) -| (stage2.north);
\draw[->,thick] (gold) -- (stage2);
\draw[->,dashed,black!55] (stage1.east) -- (stage2.west)
  node[midway,fill=white,inner sep=1pt,align=center,font=\pipelinesupport]
  {initialises weights,\\passes no output};
\draw[->,thick] (stage1) -- (answer);
\draw[->,thick] (stage2) -- (rationale);
\end{tikzpicture}
\caption[Two-stage answer and rationale pipeline]{The two-stage evaluation.
The record, question, answer choices, and any description stay the same across
the six image conditions; only the image changes. Stage~1 predicts an answer.
Stage~2 receives the gold answer and generates a rationale for it, so an image
intervention cannot also change the answer Stage~2 must support. Each input arm
has its own Stage~1 and Stage~2 adaptation. Stage~2 starts from the arm's trained
Stage~1 weights, shown by the dashed arrow; no Stage~1 answer output is passed to
Stage~2 in the primary analysis. RQ1 evaluates the answers and RQ2
evaluates the rationales.}
\label{fig:pipeline}
\end{figure}
